\section{Introduction}

The most-studied benchmarks in machine learning may be teaching models the wrong lessons. In this work, we demonstrate that models trained on high-popularity datasets exhibit generalization gaps 21 percentage points larger than those trained on less-popular alternatives from the same domain. A ResNet-18 achieving 95\% accuracy on CIFAR-10 drops to 76\% on CINIC-10, while an identical architecture trained on the less-popular SVHN maintains or even improves its accuracy on held-out data. If the benchmarks we rely on to measure progress are systematically inducing worse generalization, the field's primary progress indicators may be unreliable.

The surface-level understanding of this phenomenon points to domain shift and distribution mismatch as causes of performance degradation. However, a deeper examination reveals a more concerning pattern: popular benchmarks attract disproportionate research investment---we find 24.68 times more optimization-focused papers for high-use datasets compared to low-use alternatives ($p$=0.010). This creates what we term the \emph{benchmark co-evolution effect}, where the ML ecosystem's architecture search, hyperparameter tuning, and model design converge on dataset-specific patterns rather than generalizable features.

The gap in existing work is striking: while \citet{recht2019imagenet} documented 10--15\% accuracy drops on ImageNet with new test sets, and \citet{damour2020underspecification} provided theoretical grounding via underspecification, no systematic study has linked cross-repository dataset popularity to generalization failure. We address this gap by quantifying popularity via OpenML run-rates and arXiv paper counts, measuring generalization gaps across canonical dataset pairs, and testing the proposed co-evolution mechanism.

Our key insight is empirically grounded: high-popularity datasets (CIFAR-10) exhibit generalization gaps to held-out same-domain datasets that are 21 percentage points larger than low-popularity datasets (SVHN), despite identical training protocols and architectures. We verify the first step of the co-evolution mechanism---popular benchmarks attract significantly more optimization research (24.68$\times$, $p$=0.010)---while ruling out texture bias as the causal pathway, finding that modern architectures actually show \emph{lower} texture bias than legacy ones.

Building on these findings, we make the following contributions:

\begin{enumerate}
    \item \textbf{Empirical demonstration of the popularity-gap correlation}: We show that dataset popularity correlates with larger generalization failures across two canonical dataset pairs, with an effect size (21.21 pp) that is practically significant.

    \item \textbf{Quantification of research investment differential}: We provide the first cross-repository analysis linking dataset run-rates to arXiv optimization paper counts, finding nearly an order of magnitude difference between high-use and low-use datasets.

    \item \textbf{Mechanism falsification}: We test and refute the texture bias hypothesis at CIFAR scale, finding that skip connections in modern architectures improve rather than degrade shape-based generalization---narrowing the search space for the true causal mechanism.
\end{enumerate}

These results suggest that benchmark selection itself may be a confounding factor in model evaluation, with implications for how the research community interprets progress and how practitioners validate models before deployment.
